\begin{table}[t]
    \centering
    \resizebox{0.85\textwidth}{!}{%
    \begin{tabular}{@{}llccccc@{}}
        \toprule
        Task & Condition & Distinct conf.\ values & Conf.\ entropy & Balanced acc. & ROC-AUC & ECE \\
        \midrule
        \taskthree & \cotarm    & 10 & 1.913 & 0.666 & {\bf 0.681} & {\bf 0.083} \\
        \taskthree & \direct & {\bf 2}  & 0.644 & {\bf 0.702} & 0.597 & 0.314 \\
        \taskthree & \filler & 8  & 1.274 & 0.684 & 0.585 & 0.287 \\
        \midrule
        \taskone   & \cotarm    & 10 & 1.533 & 0.649 & 0.616 & 0.289 \\
        \taskone   & \direct & 8  & 1.189 & 0.689 & 0.676 & 0.305 \\
        \taskone   & \filler & 8  & 1.279 & {\bf 0.706} & {\bf 0.706} & 0.294 \\
        \bottomrule
    \end{tabular}
    }
    \caption{
        A confound we checked rather than assumed. AUC is computed from the model's \emph{stated}
        confidence, so confidence granularity can drive it. On \taskthree the \direct condition
        emits only two distinct confidence values, which mechanically depresses its AUC; on the
        threshold-free metric \direct is in fact better than \cotarm (0.702 vs.\ 0.666). What \cotarm does
        genuinely deliver on \taskthree is far better calibration (ECE 0.083 vs.\ 0.314)---a real
        benefit, but a different one from the hypothesised gain in discrimination.
    }
    \label{tab:calibration}
\end{table}
